\chapter{Theory Packs and Local Solver Geometries}

The repository does not attempt to solve every semantic question with one universal logic.
Instead it provides domain-specific theory packs and theory-library machinery. This is
another place where the cohomological typing viewpoint clarifies what the implementation is
doing.

A theory pack is best understood as a local geometry of solvable questions. Tensor shape,
nullability, boolean guards, witness transport, exceptions, heaps, protocols, loop
invariants, and similar concerns each induce their own families of local sections and
viability checks. The pack decides which sites are relevant, what local solver to invoke,
which restrictions are legitimate, and how trust should be upgraded when a local discharge
succeeds.

\section*{Why packs fit the site-based viewpoint}
The older monograph on sheaf-descent semantic typing already pushes hard on the idea of
site families and theory families. The hybrid and theory-pack code carries the same spirit
forward. Rather than pretending there is one maximally expressive theorem prover inside the
project, \DepPy{} distributes reasoning across packs. Each pack is responsible for a local
stalk theory.

That phrasing is not decorative. A stalk theory answers local questions about a site. A
shape theory can answer tensor-shape compatibility questions. A guard theory can answer
branch-sensitive path-viability questions. A protocol or heap theory can answer ownership
and field-availability questions. A witness transport theory can answer equality and
transport questions. The global analyzer then composes these local answers.

\section*{Trust upgrades through local discharge}
Theory packs also give a concrete answer to the question ``where do trust upgrades come
from?'' They come from successful local discharge in a domain where the relevant solver or
proof method is appropriate. A structural fact that begins as merely asserted may rise to
runtime-checked or Z3-proven status after a local pack solves it. The hybrid trust lattice
therefore does not float abstractly above the repository; it is fed by actual local
reasoning modules.

\begin{example}[Tensor theory]
A tensor-shape pack may certify that two tensors can be broadcast together, or that a
reshape preserves total size. Those are local shape facts, but once transported across
call sites and composed with code and proof layers they become part of a global contract
for a function. The same story generalizes to device compatibility, indexing safety, and
order-sensitive tensor operations.
\end{example}

\section*{Absolutizing the title carefully}
If the phrase ``absolutely everything'' is to remain honest, theory packs are part of the
answer. The repository cannot yet reason deeply about every possible Python phenomenon.
But it can repeatedly grow new local geometries and plug them into a common architecture.
The title should therefore be read as a claim about the architecture's extensibility: any
domain that can be given local sections, restrictions, and gluing conditions is in scope
for cohomological typing.
